\documentclass{article}
\usepackage{amsmath,amssymb,amsthm}
\usepackage{algorithm}
\usepackage[noend]{algpseudocode}
\usepackage{fullpage}
\usepackage{xcolor}
\newtheorem{theorem}{Theorem}
\newtheorem{lemma}{Lemma}
\newcommand{\E}{\mathbb{E}}
\newcommand{\R}{\mathbb{R}}

\begin{document}

\section*{Randomized Coordinate Descent with Importance Sampling}

\section*{Mathematical Formulation}
Let $f:\R^{n}\rightarrow\R$ be a convex differentiable function.  
For each coordinate $i\in\{1,\dots,n\}$ let $L_i>0$ be a coordinate‑wise smoothness constant, i.e.

\[
\forall \alpha\in\R,\; f(x+\alpha e_i)\le 
f(x)+\alpha\,\partial_i f(x)+\frac{L_i}{2}\alpha^{2},
\tag{1}
\]

where $e_i$ denotes the $i$‑th unit vector and $\partial_i f$ the partial derivative w.r.t. coordinate $i$.

A probability vector $p=(p_1,\dots,p_n)$ with $p_i>0$ and $\sum_{i=1}^{n}p_i=1$ is fixed beforehand.  
At iteration $k$ a coordinate $i_k$ is drawn independently from $p$.

\[
\boxed{\;
w_{k+1}=w_{k}-\eta\,\frac{1}{p_{i_k}}\,\partial_{i_k}f(w_{k})\,e_{i_k}
\;}
\tag{2}
\]

The factor $1/p_{i_k}$ guarantees that the stochastic direction is an unbiased estimator of the full gradient:
\[
\E_{i_k}\!\left[\frac{1}{p_{i_k}}\partial_{i_k}f(w_{k})e_{i_k}\right]=\nabla f(w_{k}).
\tag{3}
\]

\section*{Pseudocode}
\begin{algorithm}[H]
\caption{Randomized Coordinate Descent (RCD)}
\begin{algorithmic}[1]
\Require Initial point $w_{0}\in\R^{n}$, stepsize $\eta>0$, probability vector $p\in\Delta_{n}$
\For{$k=0,1,2,\dots,T-1$}
    \State Sample $i_k\sim p$  \Comment{draw coordinate}
    \State Compute partial derivative $g_{k} \gets \partial_{i_k}f(w_{k})$
    \State Update only coordinate $i_k$:
    \[
    w_{k+1}\gets w_{k}-\eta\,\frac{1}{p_{i_k}}\,g_{k}\,e_{i_k}
    \]
\EndFor
\State \Return $w_T$
\end{algorithmic}
\end{algorithm}

\section*{Dry Run Example}
Consider the one–dimensional quadratic
\[
f(w)=(w-3)^{2},\qquad \nabla f(w)=2(w-3).
\]
Only a single coordinate exists, thus $p_{1}=1$, $L_{1}=2$, and the update (2) reduces to ordinary gradient descent with stepsize $\eta=0.1$.

\[
\begin{array}{c|c|c|c}
k & w_{k} & \partial_{1}f(w_{k})=2(w_{k}-3) & f(w_{k}) \\ \hline
0 & 0      & -6          & 9   \\
1 & 0.6    & -4.8        & (0.6-3)^{2}=5.76 \\
2 & 1.08   & -3.84       & (1.08-3)^{2}=3.6864 \\
3 & 1.464  & -3.072      & (1.464-3)^{2}=2.3660
\end{array}
\]

The table shows the gradient computation, the coordinate update according to (2), and the resulting objective value after each iteration.

\section*{Assumptions}
\begin{enumerate}
\item \textbf{Convexity.} $f$ is convex:
\[
f(y)\ge f(x)+\langle \nabla f(x),y-x\rangle,\qquad\forall x,y\in\R^{n}.
\tag{A1}
\]
\item \textbf{Coordinate‑wise $L_i$‑smoothness.} Condition (1) holds for every $i$ with constants $L_i>0$.
\tag{A2}
\item \textbf{Probability vector.} $p_i>0$ and $\sum_{i=1}^{n}p_i=1$.
\tag{A3}
\item \textbf{Stepsize bound.} The stepsize satisfies
\[
0<\eta\le \min_{i}\frac{p_i}{L_i}.
\tag{A4}
\]
\item \textbf{Bounded optimal distance.} Let $w^{\star}\in\arg\min f$ and define
\[
R\;:=\;\max_{k\ge0}\|w_{k}-w^{\star}\|_{2}<\infty .
\tag{A5}
\]
\end{enumerate}

\section*{Main Convergence Theorem}
\begin{theorem}
\label{thm:main}
Under assumptions (A1)–(A5) and stepsize $\eta$ satisfying (A4), the iterates generated by (2) obey for every $T\ge1$
\[
\E\!\left[f(w_{T})-f(w^{\star})\right]\;\le\;
\frac{2R^{2}}{\eta\,T}.
\tag{4}
\]
Consequently the expected suboptimality decays at the rate $\mathcal{O}(1/T)$.
\end{theorem}

\section*{Theoretical Properties}
\begin{itemize}
\item \textbf{Unbiasedness.} By (3) the stochastic direction is an unbiased estimator of the full gradient, which is essential for the descent argument.
\item \textbf{Importance weighting.} The factor $1/p_i$ compensates for non‑uniform sampling; coordinates that are selected rarely receive a larger step, balancing the expected progress across all directions.
\item \textbf{Stability.} Condition (A4) guarantees that the quadratic term arising from smoothness never outweighs the linear decrease term, yielding a monotone expected descent.
\item \textbf{Comparison.} If $p_i=1/n$ for all $i$ and $L_i\equiv L$, then $\eta\le \frac{1}{nL}$ and (4) reduces to the classic $\mathcal{O}(nL R^{2}/T)$ bound for uniform random coordinate descent.
\end{itemize}

\section*{Preliminaries (Definitions and Lemmas)}
\begin{lemma}[One‑step Expected Descent]
\label{lem:descent}
For any $w\in\R^{n}$, let $i\sim p$ and define $w^{+}=w-\eta\frac{1}{p_i}\partial_i f(w)e_i$.  
If $\eta$ satisfies (A4) then
\[
\E_{i}\!\left[f(w^{+})\right]\le f(w)-\frac{\eta}{2}\,\|\nabla f(w)\|_{2}^{2}.
\tag{5}
\]
\end{lemma}
\begin{proof}
\textbf{[Step 1]} Apply coordinate‑wise smoothness (1) with $\alpha=-\eta\frac{1}{p_i}\partial_i f(w)$:
\[
f(w^{+})\le f(w)-\eta\frac{1}{p_i}(\partial_i f(w))^{2}
+\frac{L_i}{2}\eta^{2}\frac{1}{p_i^{2}}(\partial_i f(w))^{2}.
\tag{6}
\]

\textbf{[Step 2]} Take expectation w.r.t.\ $i$:
\[
\E_{i}[f(w^{+})]\le f(w)-\eta\sum_{i=1}^{n}(\partial_i f(w))^{2}
+\frac{\eta^{2}}{2}\sum_{i=1}^{n}\frac{L_i}{p_i}(\partial_i f(w))^{2}.
\tag{7}
\]

\textbf{[Step 3]} Use the stepsize bound (A4): $\displaystyle \frac{\eta L_i}{p_i}\le 1$ for all $i$. Hence
\[
\frac{\eta^{2}}{2}\sum_{i}\frac{L_i}{p_i}(\partial_i f)^{2}
\le\frac{\eta}{2}\sum_{i}(\partial_i f)^{2}.
\tag{8}
\]

\textbf{[Step 4]} Substitute (8) into (7):
\[
\E_{i}[f(w^{+})]\le f(w)-\eta\|\nabla f(w)\|_{2}^{2}
+\frac{\eta}{2}\|\nabla f(w)\|_{2}^{2}
= f(w)-\frac{\eta}{2}\|\nabla f(w)\|_{2}^{2}.
\tag{9}
\]

Thus (5) holds. ∎
\end{proof}

\begin{lemma}[Gradient lower bound via convexity]
\label{lem:gradlb}
For any $w\in\R^{n}$,
\[
\|\nabla f(w)\|_{2}^{2}\;\ge\;\frac{\bigl(f(w)-f(w^{\star})\bigr)^{2}}{R^{2}},
\tag{10}
\]
where $R$ is defined in (A5).
\end{lemma}
\begin{proof}
\textbf{[Step 5]} By convexity (A1) with $y=w^{\star}$,
\[
f(w)-f(w^{\star})\le\langle\nabla f(w),w-w^{\star}\rangle
\le\|\nabla f(w)\|_{2}\,\|w-w^{\star}\|_{2}\le\|\nabla f(w)\|_{2}\,R.
\tag{11}
\]

\textbf{[Step 6]} Rearranging (11) yields (10). ∎
\end{proof}

\section*{Main Proof (Step‑by‑Step Derivation)}
We prove Theorem~\ref{thm:main} by telescoping the one‑step descent bound.

\textbf{[Step 7]} Apply Lemma~\ref{lem:descent} with $w=w_k$ and take total expectation:
\[
\E\!\left[f(w_{k+1})\right]\le \E\!\left[f(w_{k})\right]
-\frac{\eta}{2}\,\E\!\left[\|\nabla f(w_{k})\|_{2}^{2}\right].
\tag{12}
\]

\textbf{[Step 8]} Use Lemma~\ref{lem:gradlb} inside the expectation:
\[
\E\!\left[\|\nabla f(w_{k})\|_{2}^{2}\right]
\ge \frac{1}{R^{2}}\,
\E\!\left[\bigl(f(w_{k})-f(w^{\star})\bigr)^{2}\right].
\tag{13}
\]

\textbf{[Step 9]} Insert (13) into (12):
\[
\E\!\left[f(w_{k+1})\right]\le \E\!\left[f(w_{k})\right]
-\frac{\eta}{2R^{2}}\,
\E\!\left[\bigl(f(w_{k})-f(w^{\star})\bigr)^{2}\right].
\tag{14}
\]

\textbf{[Step 10]} Denote $\Delta_k:=\E\!\left[f(w_{k})-f(w^{\star})\right]\ge0$.  
Since $\bigl(\E[X]\bigr)^{2}\le \E[X^{2}]$, we have
\[
\Delta_{k}^{2}\le \E\!\left[\bigl(f(w_{k})-f(w^{\star})\bigr)^{2}\right].
\tag{15}
\]

\textbf{[Step 11]} Using (15) in (14) gives a deterministic recursion for $\Delta_k$:
\[
\Delta_{k+1}\le \Delta_{k}-\frac{\eta}{2R^{2}}\Delta_{k}^{2}.
\tag{16}
\]

\textbf{[Step 12]} The inequality (16) is a standard discrete differential inequality.  
Define $a_k:=\Delta_k^{-1}$. Then $a_{k+1}\ge a_k+\frac{\eta}{2R^{2}}$.
Indeed,
\[
\frac{1}{\Delta_{k+1}}-\frac{1}{\Delta_{k}}
=\frac{\Delta_{k}-\Delta_{k+1}}{\Delta_{k}\Delta_{k+1}}
\ge\frac{\frac{\eta}{2R^{2}}\Delta_{k}^{2}}{\Delta_{k}^{2}}
=\frac{\eta}{2R^{2}}.
\tag{17}
\]

\textbf{[Step 13]} Summing (17) from $k=0$ to $T-1$ yields
\[
\frac{1}{\Delta_{T}}-\frac{1}{\Delta_{0}}
\ge \frac{\eta}{2R^{2}}\,T.
\tag{18}
\]

\textbf{[Step 14]} Since $\Delta_{0}=f(w_{0})-f(w^{\star})\le R^{2}L_{\max}$ is finite, we drop the negative term $1/\Delta_{0}$ to obtain the simple bound
\[
\frac{1}{\Delta_{T}}\ge \frac{\eta}{2R^{2}}\,T
\quad\Longrightarrow\quad
\Delta_{T}\le \frac{2R^{2}}{\eta\,T}.
\tag{19}
\]

\textbf{[Step 15]} Recalling $\Delta_T=\E[f(w_T)-f(w^{\star})]$ gives exactly the claim (4). ∎

\section*{Rate Derivation (With Constants)}
From (19) we have the explicit constant:
\[
\E\!\left[f(w_{T})-f(w^{\star})\right]\le
\underbrace{\frac{2R^{2}}{\eta}}_{\text{depends on }p,L_i}
\frac{1}{T}.
\]
Because $\eta\le\min_i p_i/L_i$, one may set $\eta =\frac{1}{2}\min_i \frac{p_i}{L_i}$ to obtain
\[
\E\!\left[f(w_{T})-f(w^{\star})\right]\le
\frac{4R^{2}}{T}\,\max_i\frac{L_i}{p_i}.
\tag{20}
\]

\section*{Special Cases}
\begin{itemize}
\item \textbf{Uniform sampling.} $p_i=1/n$ and $L_i\equiv L$ give $\eta\le \frac{1}{nL}$ and (20) reduces to
\[
\E[f(w_T)-f(w^{\star})]\le \frac{4nL R^{2}}{T},
\]
the classical bound for uniform random coordinate descent.
\item \textbf{Exact importance sampling.} Choose $p_i\propto L_i^{1/2}$. Then $\max_i \frac{L_i}{p_i}= \bigl(\sum_j L_j^{1/2}\bigr)^{2}$, which can be substantially smaller than $n\max_i L_i$, leading to faster practical convergence.
\end{itemize}

\section*{Discussion}
The proof hinges on two elementary ingredients: (i) coordinate‑wise smoothness to control the quadratic term, and (ii) the unbiased importance‑sampling correction $1/p_i$ to recover a full‑gradient descent‑like decrease. The resulting $\mathcal{O}(1/T)$ rate matches that of full gradient methods for convex smooth objectives, while each iteration costs only the evaluation of a single partial derivative.

\end{document}